\section*{Chapter \ref{ch:g11:polarity-and-cohesion} --- Electronegativity, Polarity and Intermolecular Forces}

\begin{solution}{exo:g11:polarity-and-cohesion:1}
\ce{H-Cl}: chlorine (3.16 against 2.20). \ce{O-H}: oxygen (3.44).
\ce{C-O}: oxygen (3.44 against 2.55).
\end{solution}

\begin{solution}{exo:g11:polarity-and-cohesion:2}
\ce{H-H}: 0, not polar. \ce{C-H}: 0.35, treated as non-polar. \ce{O-H}:
1.24, polar. \ce{N-H}: 0.84, polar. \ce{Cl-Cl}: 0, not polar. \ce{C-Cl}:
0.61, polar.
\end{solution}

\begin{solution}{exo:g11:polarity-and-cohesion:3}
Nitrogen (3.04 > 2.19) and oxygen (3.44 > 2.58): \omterm{def:g11:polarity-and-cohesion:electronegativity}{electronegativity}
decreases down a group. Nitrogen (3.04 > 2.55) and magnesium
(1.31 > 0.93): it increases from left to right along a period.
\end{solution}

\begin{solution}{exo:g11:polarity-and-cohesion:4}
Ammonia and water: each has hydrogen \omterm{def:g7:atoms-and-molecules:atom}{atoms} bonded to nitrogen or oxygen,
and \omterm{def:g10:lewis-and-shape:lone-pair}{lone pairs} on that \omterm{def:g7:atoms-and-molecules:atom}{atom}. Methane has no hydrogen on N, O or F;
hydrogen chloride has hydrogen bonded to chlorine, which is not one of
the three \omterm{def:g7:atoms-and-molecules:atom}{atoms} of the definition.
\end{solution}

\begin{solution}{exo:g11:polarity-and-cohesion:5}
Yes: \omterm{def:g11:polarity-and-cohesion:van-der-waals}{van der Waals interactions} exist between all \omterm{def:g7:atoms-and-molecules:molecule}{molecules}. But methane
\omterm{def:g7:atoms-and-molecules:molecule}{molecules} are small (10 \omterm{def:g9:inside-the-atom:nucleus}{electrons}), so these interactions are very weak
and break at very low temperature: methane boils near
\qty{-162}{\celsius}.
\end{solution}

\begin{solution}{exo:g11:polarity-and-cohesion:6}
\ce{CH2Cl2}: polar (two polar \ce{C-Cl} bonds and two nearly non-polar
\ce{C-H} bonds: no symmetry cancels them). \ce{CCl4}: non-polar (four
identical bonds at the corners of a tetrahedron). \ce{NH3}: polar
(pyramid). \ce{CO2}: non-polar (linear, opposite bonds). \ce{BF3}:
non-polar (three identical \omterm{def:g11:polarity-and-cohesion:polar-bond}{polar bonds} at \ang{120} in a plane cancel).
\end{solution}

\begin{solution}{exo:g11:polarity-and-cohesion:7}
The hydrogen of the \ce{O-H} group can be given (not those of
\ce{CH3}, bonded to carbon); the oxygen \omterm{def:g7:atoms-and-molecules:atom}{atom}, with its two \omterm{def:g10:lewis-and-shape:lone-pair}{lone pairs},
can receive. Yes: the \ce{O-H} of methanol can bond to the oxygen of a
water \omterm{def:g7:atoms-and-molecules:molecule}{molecule}, and an \ce{O-H} of water to the oxygen of methanol, for
example \ce{CH3-O-H}\,$\cdots$\,\ce{OH2}.
\end{solution}

\begin{solution}{exo:g11:polarity-and-cohesion:8}
The three \omterm{def:g7:atoms-and-molecules:molecule}{molecules} are non-polar and of the same shape, but larger and
larger (10, 18 and 36 \omterm{def:g9:inside-the-atom:nucleus}{electrons}): their \omterm{def:g11:polarity-and-cohesion:van-der-waals}{van der Waals interactions} grow,
and so do their boiling temperatures.
\end{solution}

\begin{solution}{exo:g11:polarity-and-cohesion:9}
\omterm{def:g9:periodic-table-first-look:period}{Group} 14: methane has no hydrogen bonded to N, O or F and forms no
\omterm{def:g11:polarity-and-cohesion:hydrogen-bond}{hydrogen bonds}, so it follows the trend of its group.
\end{solution}

\begin{solution}{exo:g11:polarity-and-cohesion:10}
The \omterm{def:g11:polarity-and-cohesion:van-der-waals}{van der Waals interactions}: the \omterm{def:g7:atoms-and-molecules:molecule}{molecules} grow from \ce{HCl} to
\ce{HI}, and the increase of these interactions with size outweighs the
decrease of the attraction between the \omterm{def:g11:polarity-and-cohesion:polar-bond}{partial charges}.
\end{solution}

\begin{solution}{exo:g11:polarity-and-cohesion:11}
An iodine \omterm{def:g7:atoms-and-molecules:molecule}{molecule} is much larger than a chlorine \omterm{def:g7:atoms-and-molecules:molecule}{molecule} (106
\omterm{def:g9:inside-the-atom:nucleus}{electrons} against 34): its \omterm{def:g11:polarity-and-cohesion:van-der-waals}{van der Waals interactions} are much stronger,
strong enough to hold the \omterm{def:g7:atoms-and-molecules:molecule}{molecules} in a solid at room temperature.
\end{solution}

\begin{solution}{exo:g11:polarity-and-cohesion:12}
Propane is non-polar: \omterm{def:g11:polarity-and-cohesion:van-der-waals}{van der Waals interactions} only. Methoxymethane is
polar (two polar \ce{C-O} bonds in a bent \ce{C-O-C}) but has no
hydrogen on its oxygen: no \omterm{def:g11:polarity-and-cohesion:hydrogen-bond}{hydrogen bonds} between its \omterm{def:g7:atoms-and-molecules:molecule}{molecules}. Ethanol
is polar and its \ce{O-H} groups form \omterm{def:g11:polarity-and-cohesion:hydrogen-bond}{hydrogen bonds}. The stronger the
attractions, the higher the boiling temperature: propane < methoxymethane
< ethanol.
\end{solution}

\begin{solution}{exo:g11:polarity-and-cohesion:13}
The open network of ice takes more room than the same \omterm{def:g7:atoms-and-molecules:molecule}{molecules} in the
liquid, so ice is less dense than liquid water and floats. A lake freezes
from the top: the layer of ice insulates the water underneath, which
stays liquid, and the fish survive.
\end{solution}

\begin{solution}{exo:g11:polarity-and-cohesion:14}
Solid methane (\omterm{def:g11:polarity-and-cohesion:van-der-waals}{van der Waals interactions} between small \omterm{def:g7:atoms-and-molecules:molecule}{molecules}) <
ice (\omterm{def:g11:polarity-and-cohesion:hydrogen-bond}{hydrogen bonds}) < potassium chloride (attraction between \omterm{def:g9:ions:ion}{ions}).
\end{solution}

\begin{solution}{exo:g11:polarity-and-cohesion:15}
A water \omterm{def:g7:atoms-and-molecules:molecule}{molecule} has two hydrogen \omterm{def:g7:atoms-and-molecules:atom}{atoms} to give and two \omterm{def:g10:lewis-and-shape:lone-pair}{lone pairs} to
receive with: every \omterm{def:g7:atoms-and-molecules:molecule}{molecule} can take part in four \omterm{def:g11:polarity-and-cohesion:hydrogen-bond}{hydrogen bonds}, two
given and two received, which builds a complete network. A \omterm{def:g7:atoms-and-molecules:molecule}{molecule} of
hydrogen fluoride has three \omterm{def:g10:lewis-and-shape:lone-pair}{lone pairs} but only one hydrogen \omterm{def:g7:atoms-and-molecules:atom}{atom}: it
gives a single \omterm{def:g11:polarity-and-cohesion:hydrogen-bond}{hydrogen bond}, so on average only one bond per \omterm{def:g7:atoms-and-molecules:molecule}{molecule}
can form (zig-zag chains). Water has about twice as many \omterm{def:g11:polarity-and-cohesion:hydrogen-bond}{hydrogen bonds}
per \omterm{def:g7:atoms-and-molecules:molecule}{molecule}, and boils higher.
\end{solution}

\begin{solution}{pb:g11:polarity-and-cohesion:1}
\textbf{1.} \ce{H-O-H}, \ce{H-S-H}, \ce{H-Se-H}, each central \omterm{def:g7:atoms-and-molecules:atom}{atom} with
two \omterm{def:g10:lewis-and-shape:lone-pair}{lone pairs}: O, S and Se have six \omterm{def:g10:electron-shells:valence}{valence electrons}, being in the
same group.

\textbf{2.} $M(\ce{H2O}) = \qty{18.0}{g/mol}$, $M(\ce{H2S}) =
\qty{34.1}{g/mol}$, $M(\ce{H2Se}) = \qty{81.0}{g/mol}$.

\textbf{3.} $3.44 - 2.20 = 1.24$; $2.58 - 2.20 = 0.38$;
$2.55 - 2.20 = 0.35$. Only the \ce{O-H} bond is clearly polar.

\textbf{4.} Bent (four groups, two of them \omterm{def:g10:lewis-and-shape:lone-pair}{lone pairs}). Water is polar;
hydrogen sulfide and hydrogen selenide only very slightly.

\textbf{5.} 10, 18 and 36 \omterm{def:g9:inside-the-atom:nucleus}{electrons}: hydrogen selenide has the strongest
\omterm{def:g11:polarity-and-cohesion:van-der-waals}{van der Waals interactions}.

\textbf{6.} Only water: its hydrogen \omterm{def:g7:atoms-and-molecules:atom}{atoms} are bonded to oxygen, one of
the three very electronegative \omterm{def:g7:atoms-and-molecules:atom}{atoms}; sulfur and selenium are not
electronegative enough to give their hydrogen \omterm{def:g7:atoms-and-molecules:atom}{atoms} a marked
$\delta^+$.

\textbf{7.} Four: two through its own hydrogen \omterm{def:g7:atoms-and-molecules:atom}{atoms}, two received by its
two \omterm{def:g10:lewis-and-shape:lone-pair}{lone pairs}.

\textbf{8.} \omterm{def:g11:polarity-and-cohesion:van-der-waals}{Van der Waals interactions} (including a weak attraction
between the small \omterm{def:g11:polarity-and-cohesion:polar-bond}{partial charges}).

\textbf{9.} $212.9 - 273.15 = \qty{-60.3}{\celsius}$ and
\qty{-41.3}{\celsius}.

\textbf{10.} $-41.3 - (-60.3) = \qty{19.0}{\celsius}$.

\textbf{11.} The \omterm{def:g7:atoms-and-molecules:molecule}{molecule} is larger, so its \omterm{def:g11:polarity-and-cohesion:van-der-waals}{van der Waals interactions}
are stronger.

\textbf{12.} $-60.3 - 19.0 = \qty{-79.3}{\celsius}$.

\textbf{13.} Rise $-88.1 - (-112) = \qty{23.9}{\celsius}$; prediction for
methane $-112 - 23.9 = \qty{-136}{\celsius}$, against \qty{-162}{\celsius}
in reality: the extrapolation is about \qty{26}{\celsius} too high. It is
not exact; the real trend is not a straight line.

\textbf{14.} Rise $\qty{18.7}{\celsius}$; prediction for hydrogen fluoride
$-85.1 - 18.7 = \qty{-103.8}{\celsius}$, against \qty{19.6}{\celsius}: a
gap of about \qty{123}{\celsius}.

\textbf{15.} $100 - (-79.3) = \qty{179}{\celsius}$, about
\qty{180}{\celsius}.

\textbf{16.} Rise $\qty{25.3}{\celsius}$; prediction $-87.8 - 25.3 =
\qty{-113}{\celsius}$; gap $-33.4 - (-113) = \qty{80}{\celsius}$.

\textbf{17.} Water (\qty{180}{\celsius}) > hydrogen fluoride
(\qty{123}{\celsius}) > ammonia (\qty{80}{\celsius}). Water gives two
hydrogen \omterm{def:g7:atoms-and-molecules:atom}{atoms} and receives with two \omterm{def:g10:lewis-and-shape:lone-pair}{lone pairs}: four \omterm{def:g11:polarity-and-cohesion:hydrogen-bond}{hydrogen bonds} per
\omterm{def:g7:atoms-and-molecules:molecule}{molecule}, all used. Hydrogen fluoride has one hydrogen to give, and
ammonia only one \omterm{def:g10:lewis-and-shape:lone-pair}{lone pair} to receive with: on average fewer bonds per
\omterm{def:g7:atoms-and-molecules:molecule}{molecule}, and nitrogen, less electronegative, makes weaker ones.

\textbf{18.} A gas: it would boil near \qty{-80}{\celsius}. The Earth
would have no liquid water, no oceans, no rain, and none of the life
that depends on them.

\textbf{19.} It extrapolates a straight line drawn through only two
points, and the test on \omterm{def:g9:periodic-table-first-look:period}{group} 14 shows such an extrapolation can be off
by some \qty{25}{\celsius}.

\textbf{20.} Without \omterm{def:g11:polarity-and-cohesion:hydrogen-bond}{hydrogen bonds} water would boil at about
\qty{-80}{\celsius}, some \qty{180}{\celsius} below its real boiling
point.
\end{solution}
